% General global-mode theorem (any symmetric excitatory coupling, homogeneous damping).
% Numerical check: src/theorem_general_check.py

\subsection{The confound for general excitatory coupling}\label{sec:general}

The next result states the confound without assuming that the critical mode is uniform. Its
hypothesis covers both couplings of Corollary~\ref{cor:closed} with homogeneous damping. For
diffusive coupling, $S\propto(\theta I+L)^{-1}$ has critical rate $\varepsilon=\theta$ and
$v=\one/\sqrt N$. For additive coupling, $S\propto(I-\rho\hat C)^{-1}$ has critical rate
$\varepsilon=1-\rho=\kappa$ and $v$ is the Perron vector of $C$. In both cases the eigenvectors do not
depend on the distance to the bifurcation, and
\begin{equation}\label{eq:S-eps}
  S(\varepsilon)=\varepsilon^{-1}vv^\top+R(\varepsilon),\qquad R(\varepsilon)v=0,
\end{equation}
with $v>0$, $\|v\|=1$, and $R$ real-analytic at $\varepsilon=0$ with $R(0)\succeq0$ and $\ker R(0)=\operatorname{span}\{v\}$.
A \emph{global signal} is any weighted average $g=w^\top x$ with $w$ depending analytically on
$\varepsilon$ and $w(0)^\top v\ne0$; positive weights always qualify, because $v>0$. Regressing $g$
out of the standardized signals is the GSR of Eq.~\eqref{eq:gsr}. It corresponds to
$w=D_S^{-1/2}\one\propto D_v^{-1}\one+O(\varepsilon)$, while regressing out the mean of the raw
signals corresponds to $w=\one$.

\begin{theorem}[Global-mode confound]\label{thm:confound}
Assume \eqref{eq:S-eps} and let $F(\varepsilon)=\corr S(\varepsilon)$. Write
$\tilde R=D_v^{-1}R(0)D_v^{-1}$ and $\Omega_v(i,j)=\tilde R_{ii}+\tilde R_{jj}-2\tilde R_{ij}$. As $\varepsilon\downarrow0$:
\begin{enumerate}
  \item[(a)] $F(\varepsilon)=\one\one^\top-\tfrac{\varepsilon}{2}\,\Omega_v+O(\varepsilon^2)$;
  \item[(b)] for every global signal $w$, the covariance of the residuals of $x$ after regression on $w^\top x$ satisfies
  \[
    Q(\varepsilon)=\Pi R(0)\Pi^\top+O(\varepsilon),\qquad \Pi=I-\frac{v\,w(0)^\top}{w(0)^\top v},
  \]
  so the FC after regression converges to $\corr\bigl(\Pi R(0)\Pi^\top\bigr)$, which does not depend on $\varepsilon$;
  \item[(c)] adding an independent component with loading $v$, $S\mapsto S+\sigma^2vv^\top$, is the same as replacing $\varepsilon$ by $\varepsilon_{\rm eff}$ in the singular part, $\varepsilon_{\rm eff}^{-1}=\varepsilon^{-1}+\sigma^2$, and changes $F$ by $O(\varepsilon\,\varepsilon_{\rm eff})$ relative to $F(\varepsilon_{\rm eff})$.
\end{enumerate}
The constants in (a) and (b) are controlled by $\max_i\tilde R_{ii}=\max_iR(0)_{ii}/v_i^2$, so the
expansions hold for $\varepsilon\ll\varepsilon^*:=1/\max_i\tilde R_{ii}$.
\end{theorem}

\begin{proof}
(a) Since $S_{ij}=\varepsilon^{-1}v_iv_j\bigl(1+\varepsilon\tilde R_{ij}+O(\varepsilon^2)\bigr)$ and $v_iv_j>0$,
\begin{align*}
  F_{ij}&=\frac{1+\varepsilon\tilde R_{ij}}{\sqrt{(1+\varepsilon\tilde R_{ii})(1+\varepsilon\tilde R_{jj})}}+O(\varepsilon^2)\\
  &=1-\tfrac{\varepsilon}{2}\Omega_v(i,j)+O(\varepsilon^2).
\end{align*}
(b) Write $c=v^\top w$, $r=R w$ and $\eta=\varepsilon\,w^\top r/c^2$. Then $Sw=\varepsilon^{-1}c\,v+r$ and $w^\top Sw=\varepsilon^{-1}c^2(1+\eta)$, because $v^\top r=0$. Expanding $Q=S-Sww^\top S/(w^\top Sw)$ gives the exact identity
\begin{multline*}
  Q=R-\frac{v r^\top+r v^\top}{c\,(1+\eta)}\\
  +\frac{w^\top r}{c^2(1+\eta)}\,vv^\top-\frac{\varepsilon\,rr^\top}{c^2(1+\eta)}.
\end{multline*}
The singular terms $\varepsilon^{-1}vv^\top$ cancel exactly. Letting $\varepsilon\to0$, so that $\eta\to0$, $w\to w(0)$ and $R\to R(0)$, yields $R(0)-\bigl(vr^\top+rv^\top\bigr)/c+(w^\top r/c^2)vv^\top=\Pi R(0)\Pi^\top$. The limit has positive diagonal unless some $\Pi^\top e_i$ lies in $\ker R(0)=\operatorname{span}\{v\}$, which fails for generic $w$, and $\corr$ is continuous there.
(c) $S+\sigma^2vv^\top=\varepsilon_{\rm eff}^{-1}vv^\top+R(\varepsilon)$, and $R(\varepsilon)-R(\varepsilon_{\rm eff})=O(\varepsilon)$ enters $F$ multiplied by $\varepsilon_{\rm eff}$, by (a).
\end{proof}

Three consequences deserve emphasis. First, part (b) does not require the critical mode to be close
to uniform. Any regression on a global signal with positive weights removes the divergent part of
the critical mode exactly, and after it FC no longer depends on the distance to the bifurcation at
leading order. The bound of Proposition~\ref{prop:perron} is not needed for this. Second, for
standardized signals the critical mode \emph{is} uniform at leading order: $D_S^{-1/2}v\propto\one+O(\varepsilon)$.
This is why mean FC tracks the distance to the bifurcation in both models (Sec.~\ref{sec:kappa}).
Third, Proposition~\ref{prop:perron} matters for (c). Physiological global components, whose
loading is roughly uniform, mimic a change of $\varepsilon$ to the extent that their loading overlaps
with $v$.

The window $\varepsilon^*$ can be narrow. For the Schaefer-100 group connectome,
$\varepsilon^*=2.9\times10^{-3}$ in units of $\theta$ for diffusive coupling. For additive coupling it
is $2.0\times10^{-4}$ in units of $\kappa$ for untransformed weights, whose Perron vector has
$v_{\min}/v_{\max}=0.05$, and $3.5\times10^{-3}$ after the compressive transformation
$\mathrm{SC}^{0.25}$. Outside the window the suppression by GSR is partial but already substantial.
For example, the sensitivity ratio of Proposition~\ref{prop:gsr} is 0.11 at $\theta=0.06$
(Fig.~\ref{fig:gsr}b).
